\documentclass[12pt]{article}
\usepackage[utf8]{inputenc}
\usepackage[T1]{fontenc}
\usepackage{hyperref}
\title{Wanda Language Resources for K-12 Curriculum Development}
\author{Compiled by Claude Sonnet 5 (claude-sonnet-5) for Andr\'as Kornai}
\date{}
\begin{document}
\maketitle

\section{Introduction}
Wanda (ISO 639-3: wbh; Glottocode wand1264) is a Bantu language of Tanzania, also called Ichiwanda, Iciwanda, Kiwanda, Vanda and Wandia. Speakers are concentrated in the Kamsamba ward of Momba District (Mbeya Region) and the Kipeta ward of Sumbawanga District (Rukwa Region). Most speakers know some Swahili. The searches found only general reference pages, so this report is short.

\section{Primary Educational Resources}
No primers, textbooks or school programmes for Wanda turned up in the searches.

\section{Linguistic and Reference Materials}
\begin{itemize}
\item Wikipedia entry on the Wanda language: \url{https://en.wikipedia.org/wiki/Wanda_language}
\item Wikipedia entry on the Wanda people: \url{https://en.wikipedia.org/wiki/Wanda_people}
\item Verbix endangered-languages page on Wanda: \url{https://docs.verbix.com/EndangeredLanguages/Wanda}
\item ASJP word-list entry for Wanda: \url{https://asjp.clld.org/languages/WANDA}
\end{itemize}

\section{Religious and Community Resources}
\begin{itemize}
\item Joshua Project people-group profile: \url{https://joshuaproject.net/people_groups/15851}
\item Joshua Project language page for wbh: \url{https://joshuaproject.net/languages/wbh}
\end{itemize}

\section{Digital and Audio Resources}
No audio or video resources for Wanda turned up in the searches.

\section{Curriculum Development Context}
The searches found no orthography, dictionary or school material for Wanda. Swahili is the likely language of instruction. Mother-tongue materials would need new documentation work.
\end{document}